\section{What has been established, and what has not}
\label{sec:status}

The sections above reported results of very different standing: identities verified in exact
arithmetic, numerical results with a stated accuracy, consequences of additional assumptions, and
estimates. This section sorts them. An item is called established if a machine check stands behind
it (Appendix~\ref{app:checks}); statements of interpretation deserve the reader's independent
scrutiny more than the numbers do.

\subsection{Established}

\emph{Geometry} (Sec.~\ref{sec:geometry}). The curvature and the Einstein tensor of the deflating
spacetime in closed form, for arbitrary rates and constants. The Killing vectors, 14, 15 or 20 in
number, and the completeness of the lists. That the spacetime is not a solution of Einstein's
equations without matter, with or without a cosmological constant; that with three equal rates it is
an exact solution of Einstein-Gauss-Bonnet gravity without matter, for $\Lam=3\,\gff^{2}$ and
$\aGB=-1/(4\,\gff^{2})$; and that with unequal rates it is not.
\SIG{The turning value of the ordinary time for a classical particle with momentum along a deflating
time.}{The unbounded blueshift of a classical particle that moves along a deflating coordinate, and
the incompleteness of its geodesic.}

\emph{The Dirac field} (Secs.~\ref{sec:dirac}, \ref{sec:modes}). The Dirac matrices in coordinate
components, the spin connection, the covariant constancy of the matrices, and the curvature of the
spinor derivative. The Dirac equation, the weighted field, and the statement that fields which do
not depend on $x5,x6,x7$ obey the flat equation.
\SIG{The exact solutions when one rate governs; that exactly eight of the sixteen solutions are
bounded; that these are standing waves with vanishing charge; and the law of their extinction,
Eqs.~(\ref{eq:magnitude}) to (\ref{eq:constant}).}{The exact solutions when one rate governs, and
the law of creation (\ref{eq:ncr}) with the temperature (\ref{eq:Tdefl}).} All of this is confirmed
numerically by two independent methods.

\emph{The model along the clock} (Secs.~\ref{sec:clock} to \ref{sec:cjt}). For fields that do not
depend on the extra coordinates, the model of the earlier paper on the deflating spacetime is the
same model with the couplings of the moment, exactly. The masses are constant; $\gVS$ grows as
$e^{\td/2}$; $\aVS$, $\GVS$, $\DVS$ and $C$ grow as $e^{\td}$. The modes of the Vorton-Scalar grow
with the weight only, and the potential between two fermions has the unchanged range and the
strength of the moment. The exchange amplitudes are those of flat spacetime times the square of the
coupling of the moment. In weak coupling the gap grows according to Eq.~(\ref{eq:growth}), and the
critical temperature follows it. Two fermions become bound at a definite time, at a scattering
resonance, with a binding energy that starts quadratically; the order of the other events depends on
$m_s/m$ alone. The paired state follows the deflation if and only if $\Xid$ is below about 0.2 to
0.5; before that it lags and accumulates broken pairs. The solution of the gap equation is a minimum
of the effective potential at every moment examined. Matter of fixed density loses its fermion mass
at a critical value of $C$, is self-bound for $C>3.7011$ and is more strongly bound later; the
melting temperature of the fermion mass falls; the barrier of the one-loop potential of the vacuum
moves toward the vacuum value.

\emph{Order parameter and symmetry} (Sec.~\ref{sec:order}). The potential of the order parameter
prefers no direction among the four frame axes. The internal symmetry of the sector, the
representation of the pair operators and the breaking pattern are unchanged by the deflation. Of the
six generators that belong to the pairs of the four extra axes, the deflating spacetime keeps none,
one or three as symmetries of the spacetime. A vortex line of unit winding carries four fermion zero
modes, two for each direction along the line, at every moment; its energy per length falls as
$e^{-\td}$.

\emph{Cosmology} (Sec.~\ref{sec:cosmo}), each item under the assumption named there. The equation of
the expansion (\ref{eq:friedmann}) and the fall of the inferred Hubble constant, verified by an
independent integration. The bounds of Table~\ref{tab:bounds}, and that in the model as it stands
the atomic clocks give none. The effect of the loss of mass of bound pairs. The growth of structure.
The unchanged estimates of degeneracy and of condensation in halos, the epoch of the resonance, and
the drift of the quantities.\SIG{}{ The temperature of the created quanta at the allowed rates.}

\emph{The Higgs hypothesis} (Sec.~\ref{sec:higgs}). Under H2 the electroweak scale follows the
coupling with the exponent $\nEW$, and atomic clocks bound the rate of deflation by
Eq.~(\ref{eq:clockbound}). The exponent is given in closed form for the scale-invariant ladder
equation.

\subsection{Not established}

\emph{The source of the deflation.} The metric is prescribed. For equal rates it solves the
equations of Einstein-Gauss-Bonnet gravity without matter; whether that solution is stable has not
been examined. For unequal rates no field equation has been found that the metric solves, and the
stress tensor that Einstein's equations would require has not been identified with any matter. All
results on fields are results on a given background.

\emph{The theory of gravity and the assumption A2.} If gravity is of the Einstein-Gauss-Bonnet kind,
the relation (\ref{eq:GN}) between Newton's constant and the volume of the deflating coordinates,
and with it every result of Sec.~\ref{sec:cosmo} that is marked A2, would have to be derived anew.
Only the time component of Einstein's equations was used; whether a source exists that keeps three
rates constant while ordinary space expands has not been examined.

%%<1644>
\emph{The exclusion of the growing solutions.} The eight growing solutions of each mode that depends
on a deflating time are discarded because they are unbounded. That is a condition of admissibility;
it is not derived from a principle. The theory has no positive conserved charge for such modes, and
a consistent quantization exists only for the sector that does not depend on the additional times.
%%</1644>
%%<1671>
\emph{The reaction of the created particles on the metric.} The particles created by the deflation
gain energy without bound. Their stress must eventually change the metric; this has not been
computed. Nor do the coordinates cover the whole future of particles that move along deflating
coordinates; what lies beyond is outside the hypothesis.

\emph{The restriction to the sector.} That the fields do not depend on $x5,x6,x7$ is assumed, as in
the earlier paper. The deflation supports the assumption only if the deflating coordinates are
periodic (Sec.~\ref{sec:modes}).
%%</1671>

\emph{Unequal rates with simultaneous crossings.} No closed form is known for a mode that depends on
two deflating coordinates whose crossings coincide; the results are numerical.

\emph{The hidden coordinate.} That the fields do not depend on $x0$ is assumed, as in the earlier
paper. Nothing in the deflation explains it.

\emph{The end of the history.} The couplings grow without bound. The four-dimensional description
with the exchange of one boson holds only while $\aVS$ is not large, and the eight-dimensional theory
itself is an effective theory, with couplings of negative mass dimension, that must fail when the
proper size of the deflating coordinates reaches its scale of validity. What happens then is outside
the model. Every history of Secs.~\ref{sec:gap} to \ref{sec:cjt} ends in a regime that the two-loop
truncation does not control; the numbers at the last moments of each table indicate the trend and
are not predictions.

\emph{The energy.} The energy of matter is not conserved while the couplings change. Where it goes,
into whatever drives the deflation, is not described.

\emph{Dynamics.} Sections~\ref{sec:gap}, \ref{sec:pair} and \ref{sec:cjt} give sequences of
equilibrium states. Only for the paired state has the ability to follow been examined, in a model
without dissipation; with dissipation the broken pairs would relax, and their lifetime is not known.
The time-dependent two-body problem at the moment of binding, which decides how much of the
population ends in the bound state, has not been solved. How a gas reaches the self-bound state, and
whether the abnormal state forms by nucleation, are open. The equilibrium of a loop of vortex line
with its current, and its stability, have not been computed.

\emph{The vacuum beyond one loop.} The statements about the vacuum use the one-loop potential
outside its range of reliability and draw only the conclusion stated in Sec.~\ref{sec:cjt}.

\emph{Cosmology.} The spectra of the microwave background with a Newton's constant that depends on
time have not been computed (Sec.~\ref{sec:cosmo}). All bounds that reach back to nucleosynthesis
assume a constant rate; a rate that was large early and is small today is neither excluded nor
described. The consequences of the epoch of the resonance for the formation of halos have not been
simulated.

\emph{The Higgs hypothesis.} Everything that the earlier paper lists as not established: that a
condensate forms in the vacuum, that \SIG{any symmetry of the model}{the rotation group of the extra
dimensions} is the electroweak gauge group, the scale and the mass of the scalar particle, and the
masses of ordinary fermions. The exponent $\nEW$ has not been computed for the model itself; the
statement that it is of order one or larger rests on the scale-invariant case and on the general
fact that the scale of the condensate vanishes at the critical coupling.

\section{Conclusions}
\label{sec:concl}

We asked what becomes of a 16-component Dirac field, and of the superfluid of its pairs, when three
of the eight coordinates of a spacetime of signature \SIG{$4+4$}{$7+1$} deflate exponentially.

The geometry is simple and rigid. The deflating spacetime is homogeneous, it has 14 to 20 symmetries
instead of 36, and it is not a solution of Einstein's equations without matter. With three equal
rates it is an exact solution of Einstein-Gauss-Bonnet gravity without matter, for one relation
between the rate and the two constants of that theory.

The Dirac field separates into two parts. Fields that do not depend on the deflating coordinates
obey, after a change of normalization, exactly the flat equations.
\SIG{Fields that do depend on them are extinguished: every admissible mode is a standing wave that
dies out as the exponential of an exponential of the time. The sector in which the $4+4$ theory is
not unitary is removed by the deflation, although the condition that selects the admissible modes is
imposed and not derived.}{Fields that do depend on them are populated by the deflation with a
Fermi-Dirac spectrum at the temperature of the rate divided by $2\pi$, and their quanta then gain
energy without bound.}

The superfluid of hypothesis H1 exists on the deflating spacetime at every moment, with the
couplings of that moment, and it acquires a history. The masses do not change; the coupling grows
with the inverse volume of the deflating coordinates; the gap grows much faster than the coupling,
in weak coupling as the exponential of an exponential; two fermions become bound at a scattering
resonance; matter loses its fermion mass and
becomes self-bound; and the paired state can follow the deflation only once its gap exceeds a few
times the rate divided by the pairing strength. The direction of the order parameter is not selected
by the deflation.

For hypotheses H2 and H3 the deflation brings no relief. \SIG{The Higgs hypothesis had no geometric
basis on the flat $4+4$ spacetime and has none here.}{The rotations that gave the Higgs hypothesis
its one structural support on the flat $7+1$ spacetime are no longer symmetries of the spacetime.}
If the superfluid were the Higgs field, the electroweak scale would follow the growing coupling, and
atomic clocks would then limit the rate of deflation to about a millionth of the Hubble rate. If
gravity is the Einstein theory in eight dimensions, lunar laser ranging limits it to three
thousandths of the Hubble rate, and the Hubble constant inferred from the microwave background
moves away from the local measurement. In either case nothing of the history just described can
have happened within the age of the universe at a constant rate.

These results are quantitative, and they come with the programs that produced them. They mark what
a successful theory of this kind would have to supply: a source or a theory of gravity that sustains
the deflation, a rate that was large early and is small today, and a reason why the constants of
ordinary matter do not follow the volume of the deflating coordinates.

\section*{Data availability}

The source code of all programs, the data files from which every figure and every computed number of
this paper is generated, the instructions for compiling and running the programs, the supplements
that derive each result step by step, and the documents that record the provenance of each result
are openly available \cite{DeflationRepo}.

